\chapter{Wymiar topologiczny}
\label{ch:topog23-wymiar}

\noindent\textbf{Cel rozdziału.}
W przestrzeni euklidesowej liczba współrzędnych sugeruje
pojęcie wymiaru. Dla dowolnych przestrzeni topologicznych
nie mamy jednak wyróżnionych współrzędnych.
Wprowadzimy trzy niezmienniki oparte wyłącznie na topologii:
wymiar pokryciowy Lebesgue'a, mały wymiar indukcyjny
oraz duży wymiar indukcyjny. Wyjaśnimy ich sens,
podamy ścisłe charakterystyki wymiaru zerowego
i przeprowadzimy samodzielne dowody wymiaru odcinka
oraz górnego oszacowania wymiaru zwartych podzbiorów
przestrzeni euklidesowych. Pokażemy także rolę
nerwu pokrycia i map do kompleksów symplicjalnych.
W głównych twierdzeniach o wymiarze pokryciowym
pracujemy z przestrzeniami zwartymi Hausdorffa,
aby wszystkie argumenty o pokryciach miały jawnie
podane założenia.

\section{Rząd pokrycia i wymiar Lebesgue'a}
\label{sec:topog23-cover}

\begin{definicja}[Pokrycia, rozdrobnienia i rząd]
\label{def:topog23-order}
Niech \(X\) będzie przestrzenią topologiczną.
Rodzinę \(\mathcal U=(U_i)_{i\in I}\) nazywamy
\emph{pokryciem} \(X\), jeśli
\(X=\bigcup_{i\in I}U_i\).
Pokrycie \(\mathcal V=(V_j)_{j\in J}\)
jest \emph{rozdrobnieniem} \(\mathcal U\),
jeżeli dla każdego \(j\in J\) istnieje
\(i\in I\) takie, że \(V_j\subseteq U_i\).
Rodziny traktujemy jako indeksowane,
więc dwa wystąpienia tego samego zbioru
z różnymi indeksami liczymy osobno.

\emph{Rząd pokrycia} \(\mathcal U\)
jest równy co najwyżej \(m\in\mathbb N\),
jeżeli każdy punkt \(x\in X\)
należy do co najwyżej \(m\) zbiorów \(U_i\).
Inaczej: dla każdych różnych indeksów
\(i_0,\ldots,i_m\in I\) zachodzi
\(U_{i_0}\cap\cdots\cap U_{i_m}=\varnothing\).
\end{definicja}

\begin{definicja}[Wymiar pokryciowy Lebesgue'a]
\label{def:topog23-dim}
Niech \(X\) będzie zwartą przestrzenią Hausdorffa.
Przyjmujemy \(\dim\varnothing=-1\).
Dla niepustego \(X\) i \(n\in\{0,1,2,\ldots\}\)
piszemy \(\dim X\le n\), jeżeli
każde skończone pokrycie otwarte \(X\)
ma skończone otwarte rozdrobnienie
rzędu co najwyżej \(n+1\).
Liczbą \(\dim X\) nazywamy najmniejsze
takie \(n\), o ile istnieje;
w przeciwnym razie piszemy \(\dim X=\infty\).
\end{definicja}

\begin{uwaga}[Zakres definicji]
\label{rem:topog23-normality}
Dla przestrzeni niezwartych używa się również
ogólnej definicji wymiaru pokryciowego,
z odpowiednimi warunkami na rozdrobnienia.
Tutaj ograniczenie do zwartych przestrzeni Hausdorffa
zapewnia równoważność sformułowań przez
pokrycia skończone i dowolne pokrycia otwarte:
z dowolnego pokrycia wybiera się skończone
podpokrycie, które następnie rozdrabnia się
do żądanego rzędu.
Nie należy utożsamiać tego pojęcia
z wymiarem Hausdorffa metrycznej przestrzeni.
\end{uwaga}

\begin{stwierdzenie}[Niezmienniczość homeomorficzna]
\label{prop:topog23-homeomorphism}
Niech \(X,Y\) będą zwartymi przestrzeniami Hausdorffa
i \(h\colon X\to Y\) homeomorfizmem.
Wówczas \(\dim X=\dim Y\).
\end{stwierdzenie}
\begin{proof}
Załóżmy, że \(\dim X\le n\).
Dla skończonego otwartego pokrycia
\((U_i)\) przestrzeni \(Y\)
przeciwobrazy \(h^{-1}(U_i)\)
tworzą skończone otwarte pokrycie \(X\).
Wybierzmy jego skończone otwarte rozdrobnienie
\((V_j)\) rzędu co najwyżej \(n+1\).
Rodzina \((h(V_j))\) jest otwartym
pokryciem \(Y\), rozdrabnia \((U_i)\),
a bijektywność \(h\) zachowuje
liczbę zbiorów zawierających punkt.
Stąd \(\dim Y\le n\).
Stosujemy ten sam argument do \(h^{-1}\).
W przypadku wymiaru nieskończonego
równość wynika z obu implikacji.
Przypadek pusty jest oczywisty.
\end{proof}

\begin{lemat}[Zamknięte zmniejszenie skończonego pokrycia]
\label{lem:topog23-closed-shrinking}
Niech \(X\) będzie zwartą przestrzenią Hausdorffa,
a \(U_1,\ldots,U_m\subseteq X\)
skończonym pokryciem otwartym,
gdzie \(m\ge1\).
Istnieją zbiory domknięte \(F_i\subseteq U_i\)
dla \(1\le i\le m\), takie że
\(X=\bigcup_{i=1}^mF_i\).
\end{lemat}
\begin{proof}
Przestrzeń \(X\) jest normalna
na mocy Twierdzenia~\ref{thm:topog-zwartosc-hausdorff-normalna}.
Dla każdego \(x\in X\) wybieramy \(i(x)\)
z \(x\in U_{i(x)}\).
Regularność pozwala wybrać
otwarty \(W_x\ni x\) spełniający
\(\overline{W_x}^{\,X}\subseteq U_{i(x)}\).
Ze zwartości \(X\) istnieją
\(x_1,\ldots,x_k\in X\) takie, że
\(X=\bigcup_{j=1}^k W_{x_j}\).
Dla \(1\le i\le m\) definiujemy
\[
 F_i=\bigcup_{\{j:i(x_j)=i\}}
           \overline{W_{x_j}}^{\,X},
\]
przyjmując pusty zbiór, gdy wskazana rodzina
indeksów jest pusta.
Każdy \(F_i\) jest skończoną sumą
zbiorów domkniętych, więc jest domknięty.
Z konstrukcji \(F_i\subseteq U_i\),
a suma wszystkich \(F_i\) zawiera
sumę zbiorów \(W_{x_j}\), czyli całe \(X\).
\end{proof}

\begin{twierdzenie}[Dziedziczenie przez zbiory domknięte]
\label{thm:topog23-closed-monotone}
Niech \(X\) będzie zwartą przestrzenią Hausdorffa,
a \(F\subseteq X\) zbiorem domkniętym.
Wówczas \(\dim F\le\dim X\)
w znaczeniu: jeśli \(\dim X\le n\)
dla pewnego \(n\ge0\), to \(\dim F\le n\).
\end{twierdzenie}
\begin{proof}
Gdy \(F=\varnothing\), teza wynika z definicji.
Dla niepustego \(F\) weźmy dowolne
skończone otwarte pokrycie
\(U_1,\ldots,U_m\) przestrzeni \(F\).
Istnieją otwarte \(W_i\subseteq X\)
takie, że \(U_i=W_i\cap F\).
Rodzina \(W_1,\ldots,W_m,X\setminus F\)
jest skończonym otwartym pokryciem \(X\).
Wybierzmy otwarte rozdrobnienie
\((V_j)_{j=1}^k\) rzędu co najwyżej \(n+1\).
Zbiory \(V_j\cap F\) niepuste,
po odrzuceniu tych \(V_j\),
które rozdrabniają \(X\setminus F\),
tworzą otwarte pokrycie \(F\).
Każdy z nich zawiera się w jednym z \(U_i\).
Ich rząd nie przekracza rzędu \((V_j)\),
więc \(\dim F\le n\).
\end{proof}

\section{Wymiar zero i bazy zbiorów otwarto-domkniętych}
\label{sec:topog23-zero}

\begin{twierdzenie}[Charakteryzacja wymiaru zerowego]
\label{thm:topog23-zero-clopen}
Niech \(X\) będzie niepustą zwartą przestrzenią Hausdorffa.
Następujące warunki są równoważne:
\begin{enumerate}[label=\textup{(\roman*)}]
\item \(\dim X=0\);
\item każdy punkt ma bazę otoczeń
złożoną ze zbiorów jednocześnie otwartych
i domkniętych w \(X\);
\item każde skończone otwarte pokrycie \(X\)
ma skończone rozdrobnienie będące
pokryciem przez parami rozłączne
zbiory otwarto-domknięte.
\end{enumerate}
\end{twierdzenie}
\begin{proof}
Załóżmy \textup{(i)}.
Niech \(x\in U\subseteq X\), gdzie \(U\) jest otwarty.
Zwartość Hausdorffa implikuje aksjomat \(T_1\),
więc \(\{U,X\setminus\{x\}\}\)
jest skończonym otwartym pokryciem \(X\).
Istnieje skończone otwarte rozdrobnienie
rzędu co najwyżej \(1\).
Ponieważ rozdrobnienie pokrywa \(X\),
jego niepuste elementy są parami rozłączne.
Jest wśród nich dokładnie jeden \(V\ni x\).
Zbiór \(V\) musi zawierać się w \(U\),
gdyż \(X\setminus\{x\}\) nie zawiera \(x\).
Dopełnienie \(V\) jest sumą pozostałych
elementów otwartego pokrycia,
zatem jest otwarte.
Wobec tego \(V\) jest otwarto-domknięty
i \(x\in V\subseteq U\).
Dowiedliśmy \textup{(ii)}.

Załóżmy \textup{(ii)} i rozważmy
skończone otwarte pokrycie \((U_i)_{i=1}^m\).
Dla każdego \(x\in X\) możemy wybrać
otwarto-domknięty \(C_x\) taki, że
\(x\in C_x\subseteq U_{i(x)}\).
Ze zwartości \(X\) wybieramy skończone
podpokrycie \(C_1,\ldots,C_k\).
Definiujemy
\[
 D_1=C_1,\qquad
 D_j=C_j\setminus\bigcup_{l<j}C_l
 \quad (2\le j\le k).
\]
Zbiory \(D_j\) są otwarto-domknięte,
parami rozłączne i pokrywają \(X\);
każdy zawiera się w odpowiednim \(C_j\),
a więc rozdrabnia pierwotne pokrycie.
Możemy pominąć wyrazy puste.
Dostajemy \textup{(iii)}.

Warunek \textup{(iii)} jest wprost
warunkiem \(\dim X\le0\).
Ponieważ \(X\neq\varnothing\),
otrzymujemy \(\dim X=0\).
\end{proof}

\begin{stwierdzenie}[Przestrzenie skończone i dyskretne]
\label{prop:topog23-finite-zero}
Każda niepusta skończona przestrzeń dyskretna
ma wymiar pokryciowy \(0\).
\end{stwierdzenie}
\begin{proof}
Przestrzeń skończona dyskretna jest zwarta
i Hausdorffa. Jej wszystkie singletony
są otwarto-domknięte, a zatem tworzą bazę.
Twierdzenie~\ref{thm:topog23-zero-clopen}
daje tezę.
\end{proof}

\begin{stwierdzenie}[Iloczyny przestrzeni wymiaru zero]
\label{prop:topog23-zero-products}
Niech \((X_i)_{i\in I}\) będzie dowolną
rodziną niepustych zwartych przestrzeni
Hausdorffa o wymiarze pokryciowym \(0\).
Wówczas iloczyn
\(X=\prod_{i\in I}X_i\)
z topologią produktową ma wymiar \(0\).
Dotyczy to także \(I=\varnothing\),
gdy \(X\) jest jednoelementowy.
\end{stwierdzenie}
\begin{proof}
Twierdzenie Tychonoffa
\ref{thm:topog-tychonoff-iloczyny}
zapewnia zwartość \(X\),
a własność Hausdorffa przechodzi na iloczyny.
Ponieważ każdy \(X_i\) ma bazę
zbiorów otwarto-domkniętych,
dla każdego punktu \(x\in X\)
i podstawowego otoczenia
\[
 O=\bigcap_{j\in J}p_j^{-1}(U_j),
 \qquad J\subseteq I\text{ skończony},
\]
możemy wybrać otwarto-domknięte
\(C_j\subseteq U_j\) z \(x_j\in C_j\).
Wówczas
\(C=\bigcap_{j\in J}p_j^{-1}(C_j)\)
jest zbiorem otwarto-domkniętym
z \(x\in C\subseteq O\).
Zbiory te tworzą bazę iloczynu.
Na mocy Twierdzenia~\ref{thm:topog23-zero-clopen}
wymiar \(X\) jest zerowy.
W przypadku \(I=\varnothing\)
mamy przestrzeń jednoelementową,
a więc również wymiar zero.
\end{proof}

\begin{przyklad}[Zbiór Cantora]
\label{ex:topog23-cantor}
Niech
\[
 C=\left\{\sum_{j=1}^{\infty}
          \frac{2\varepsilon_j}{3^j}:
          \varepsilon_j\in\{0,1\}\right\}
 \subseteq[0,1]
\]
będzie klasycznym zbiorem Cantora.
Dla każdej skończonej listy cyfr
\(\varepsilon_1,\ldots,\varepsilon_m\)
zbiór punktów \(C\) o wskazanych
pierwszych \(m\) cyfrach jest
otwarto-domknięty w \(C\).
Istotnie, jest obrazem podstawowego
zbioru cylindrycznego w
\(\{0,1\}^{\mathbb N}\).
Odwzorowanie kodujące cyfry jest ciągłe,
gdyż ogon szeregu ma moduł co najwyżej \(3^{-m}\);
jest też różnowartościowe:
przy pierwszej różnicy na miejscu \(m\)
różnica odpowiadających sum ma moduł
co najmniej
\(2\cdot3^{-m}-\sum_{j>m}2\cdot3^{-j}
=3^{-m}>0\).
Jest zatem homeomorfizmem zwartej przestrzeni
\(\{0,1\}^{\mathbb N}\) na \(C\).
Podstawowe cylindry są otwarto-domknięte,
stąd \(\dim C=0\).
Równoważnie, można zastosować
Stwierdzenie~\ref{prop:topog23-zero-products}
i niezmienniczość homeomorficzną.
\end{przyklad}

\begin{twierdzenie}[Niezerowy wymiar przestrzeni spójnej]
\label{thm:topog23-connected-positive}
Niech \(X\) będzie zwartą przestrzenią Hausdorffa
zawierającą przynajmniej dwa punkty.
Jeśli \(X\) jest spójna, to \(\dim X\ge1\).
\end{twierdzenie}
\begin{proof}
Wybierzmy różne \(a,b\in X\).
Zbiory \(X\setminus\{a\}\) i
\(X\setminus\{b\}\) są otwarte i właściwe,
a ich suma pokrywa \(X\).
Gdyby \(\dim X=0\), na mocy
Twierdzenia~\ref{thm:topog23-zero-clopen}
pokrycie miałoby rozdrobnienie
przez parami rozłączne
zbiory otwarto-domknięte.
Ponieważ \(X\) jest spójna,
każdy taki zbiór jest pusty albo równy \(X\).
Niepuste rozdrobnienie musiałoby
więc składać się z samego \(X\),
ale \(X\) nie jest zawarty w żadnym
z dwóch wyjściowych zbiorów.
Otrzymujemy sprzeczność.
\end{proof}

\section{Wymiary indukcyjne}
\label{sec:topog23-inductive}

\begin{definicja}[Mały wymiar indukcyjny]
\label{def:topog23-ind}
Niech \(X\) będzie zwartą przestrzenią Hausdorffa.
Definiujemy rekurencyjnie
\(\operatorname{ind}\varnothing=-1\).
Dla \(n\ge0\) piszemy
\(\operatorname{ind}X\le n\), jeśli
dla każdego \(x\in X\) i każdego otwartego
otoczenia \(U\ni x\) istnieje otwarty zbiór \(V\)
spełniający
\[
 x\in V\subseteq\overline V^{\,X}\subseteq U,
 \qquad
 \operatorname{ind}(\partial_XV)\le n-1,
\]
gdzie
\(\partial_XV=\overline V^{\,X}\setminus V\).
Mały wymiar indukcyjny
\(\operatorname{ind}X\) jest najmniejszym
takim \(n\), a w przypadku braku
jakiegokolwiek \(n\) wynosi \(\infty\).
\end{definicja}

\begin{definicja}[Duży wymiar indukcyjny]
\label{def:topog23-Ind}
Niech \(X\) będzie zwartą przestrzenią Hausdorffa.
Przyjmujemy \(\operatorname{Ind}\varnothing=-1\).
Dla \(n\ge0\) piszemy
\(\operatorname{Ind}X\le n\), jeżeli
dla każdej pary \(F\subseteq U\subseteq X\),
gdzie \(F\) jest domknięty, a \(U\) otwarty,
istnieje otwarty \(V\subseteq X\) taki, że
\[
 F\subseteq V\subseteq\overline V^{\,X}\subseteq U,
 \qquad
 \operatorname{Ind}(\partial_XV)\le n-1.
\]
Definiujemy \(\operatorname{Ind}X\)
jako najmniejsze odpowiednie \(n\)
lub \(\infty\), jeśli takie \(n\) nie istnieje.
\end{definicja}

\begin{uwaga}[Znaczenie warunku o brzegu]
\label{rem:topog23-boundary-ind}
Wymiar indukcyjny opisuje nie tylko
liczbę zbiorów spotykających się w punkcie,
lecz możliwość wyboru otoczeń,
których brzeg ma mniejszy wymiar.
Dla \(n=0\) warunek
\(\operatorname{ind}(\partial_XV)\le-1\)
oznacza \(\partial_XV=\varnothing\).
Wtedy \(V\) jest otwarto-domknięty.
Wymiar mały dotyczy otoczeń punktów;
wymiar duży --- oddzielania
całych zbiorów domkniętych od dopełnienia
otwartego otoczenia.
\end{uwaga}

\begin{twierdzenie}[Zgodność trzech wymiarów w wymiarze zero]
\label{thm:topog23-three-zero}
Niech \(X\) będzie niepustą zwartą
przestrzenią Hausdorffa.
Wówczas następujące warunki są równoważne:
\[
 \dim X=0,\qquad
 \operatorname{ind}X=0,\qquad
 \operatorname{Ind}X=0.
\]
\end{twierdzenie}
\begin{proof}
Z definicji małego wymiaru indukcyjnego
warunek \(\operatorname{ind}X\le0\)
jest równoważny istnieniu bazy
otwarto-domkniętych otoczeń każdego punktu.
Istotnie, w takim przypadku
\(\partial_XV=\varnothing\), a ponieważ
\(V\) jest otwarty, oznacza to,
że \(V\) jest także domknięty.
Odwrotnie, otwarto-domknięte \(V\)
ma pusty brzeg.
Twierdzenie~\ref{thm:topog23-zero-clopen}
daje równoważność
\(\dim X=0\Longleftrightarrow
 \operatorname{ind}X=0\).

Jeśli \(\operatorname{Ind}X=0\),
stosujemy definicję dużego wymiaru
do domkniętego singletonu \(F=\{x\}\)
i dowolnego otwartego \(U\ni x\).
Otrzymujemy otwarto-domknięty zbiór \(V\)
z \(x\in V\subseteq U\);
stąd \(\operatorname{ind}X=0\).

Załóżmy wreszcie \(\dim X=0\).
Niech \(F\subseteq U\subseteq X\),
gdzie \(F\) jest domknięty, a \(U\) otwarty.
Zwartość \(F\) oraz baza zbiorów
otwarto-domkniętych pozwalają wybrać
skończenie wiele takich zbiorów
\(C_1,\ldots,C_k\subseteq U\),
które pokrywają \(F\).
Ich suma \(V=\bigcup_{j=1}^k C_j\)
jest otwarto-domknięta,
a \(F\subseteq V\subseteq U\).
W szczególności \(\partial_XV=\varnothing\),
więc \(\operatorname{Ind}X\le0\).
Niepustość \(X\) wyklucza wymiar \(-1\),
co kończy dowód.
\end{proof}

\section{Lemat o liczbie Lebesgue'a i wymiar odcinka}
\label{sec:topog23-interval}

\begin{lemat}[Liczba Lebesgue'a pokrycia]
\label{lem:topog23-lebesgue-number}
Niech \((K,d)\) będzie niepustą
zwartą przestrzenią metryczną,
a \((U_i)_{i=1}^m\) jej skończonym
otwartym pokryciem.
Istnieje \(\delta>0\) takie,
że każdy niepusty podzbiór \(A\subseteq K\)
o średnicy mniejszej niż \(\delta\)
jest zawarty w pewnym \(U_i\).
\end{lemat}
\begin{proof}
Dla każdego \(x\in K\) wybieramy \(i(x)\)
oraz \(\varepsilon_x>0\) tak, że
\(B_d(x,2\varepsilon_x)\subseteq U_{i(x)}\).
Kule \(B_d(x,\varepsilon_x)\) pokrywają \(K\).
Ze zwartości istnieją punkty
\(x_1,\ldots,x_k\) takie, że
\(K=\bigcup_{j=1}^k B_d(x_j,\varepsilon_{x_j})\).
Przyjmijmy
\(\delta=\min_{1\le j\le k}\varepsilon_{x_j}>0\).

Jeśli \(A\neq\varnothing\) i
\(\operatorname{diam} A<\delta\),
wybierzmy \(a\in A\) i indeks \(j\)
z \(a\in B_d(x_j,\varepsilon_{x_j})\).
Dla każdego \(y\in A\)
\[
 d(y,x_j)\le d(y,a)+d(a,x_j)
 <\delta+\varepsilon_{x_j}
 \le2\varepsilon_{x_j}.
\]
Zatem \(A\subseteq U_{i(x_j)}\).
\end{proof}

\begin{twierdzenie}[Wymiar zwartych podzbiorów prostej]
\label{thm:topog23-real-upper}
Niech \(K\subseteq\mathbb R\) będzie zwartym
podzbiorem prostej z jej zwykłą topologią.
Wówczas \(\dim K\le1\),
przy konwencji \(\dim\varnothing=-1\).
\end{twierdzenie}
\begin{proof}
Dla pustego \(K\) teza jest oczywista.
Załóżmy \(K\neq\varnothing\)
i weźmy skończone otwarte pokrycie \((U_i)\)
przestrzeni \(K\).
Niech \(\delta>0\) będzie liczbą
z Lematu~\ref{lem:topog23-lebesgue-number}
dla metryki euklidesowej.
Wybierzmy \(h>0\) z \(2h<\delta\).
Dla każdego \(k\in\mathbb Z\) definiujemy
\[
 V_k=K\cap(kh-h,kh+h).
\]
Zbiory \(V_k\) są otwarte w \(K\) i pokrywają \(K\):
każda liczba rzeczywista leży w odległości
co najwyżej \(h/2\) od pewnej wielokrotności \(h\).
Ponieważ \(K\) jest ograniczony,
tylko skończenie wiele \(V_k\)
jest niepustych.

Średnica każdego \(V_k\)
jest mniejsza lub równa \(2h<\delta\),
więc niepuste \(V_k\) zawiera się w pewnym \(U_i\).
Ponadto każdy punkt należy
do co najwyżej dwóch zbiorów \(V_k\):
gdyby \(x\) należał do \(V_{k_1}\)
i \(V_{k_3}\) dla \(k_3\ge k_1+2\),
musielibyśmy mieć jednocześnie
\(|x-k_1h|<h\) i
\(|x-k_3h|<h\), co przez nierówność trójkąta
dawałoby \((k_3-k_1)h<2h\),
sprzeczność.
Pokrycie \((V_k)\) jest więc
skończonym otwartym rozdrobnieniem
rzędu co najwyżej \(2\).
To dowodzi \(\dim K\le1\).
\end{proof}

\begin{twierdzenie}[Wymiar odcinka]
\label{thm:topog23-interval}
Dla dowolnych \(a,b\in\mathbb R\) z \(a<b\)
zachodzi
\[
 \dim[a,b]=1.
\]
\end{twierdzenie}
\begin{proof}
Twierdzenie~\ref{thm:topog23-real-upper}
daje \(\dim[a,b]\le1\).
Odcinek jest zwartą przestrzenią Hausdorffa,
spójną i zawiera więcej niż jeden punkt.
Twierdzenie~\ref{thm:topog23-connected-positive}
wyklucza wymiar zero.
Otrzymujemy równość.
\end{proof}

\begin{wniosek}[Zwarte spójne podzbiory prostej]
\label{cor:topog23-connected-real}
Niech \(K\subseteq\mathbb R\) będzie niepusty,
zwarty i spójny.
Jeżeli \(K\) ma jeden punkt,
to \(\dim K=0\).
Jeżeli \(K\) zawiera przynajmniej dwa punkty,
to \(\dim K=1\).
\end{wniosek}
\begin{proof}
Jednopunktowa przestrzeń jest dyskretna,
więc ma wymiar zero.
Jeśli \(K\) ma co najmniej dwa punkty,
jego spójność i Twierdzenie
\ref{thm:topog23-connected-positive}
dają \(\dim K\ge1\).
Z Twierdzenia~\ref{thm:topog23-real-upper}
mamy \(\dim K\le1\).
\end{proof}

\begin{figure}[htbp]
\centering
\begin{tikzpicture}[x=1.02cm,y=.7cm]
 \draw[->] (-.5,0)--(7.4,0);
 \foreach \k in {0,...,6}{
  \draw[very thin,gray] (\k,-.12)--(\k,.14);
  \draw[blue!65!black,very thick] (\k-.87,.55)--(\k+.87,.55);
  \fill[blue!65!black] (\k-.87,.55) circle (1.4pt);
  \fill[blue!65!black] (\k+.87,.55) circle (1.4pt);
 }
 \node[anchor=west,font=\small] at (0,-.56)
 {Punkty prostej należą do najwyżej dwóch kolejnych otoczeń.};
\end{tikzpicture}
\caption{Idea rozdrobnienia pokrycia zwartego podzbioru
prostej do rzędu co najwyżej \(2\).}
\label{fig:topog23-line-cover}
\end{figure}

\section{Nerw pokrycia i odwzorowanie barycentryczne}
\label{sec:topog23-nerve}

\begin{definicja}[Nerw skończonego pokrycia]
\label{def:topog23-nerve}
Niech \(\mathcal U=(U_i)_{i=1}^{m}\)
będzie skończonym otwartym pokryciem \(X\).
\emph{Nerwem} \(\mathcal U\) nazywamy
abstrakcyjny kompleks symplicjalny,
którego wierzchołki odpowiadają indeksom
\(1,\ldots,m\), a niepusty zbiór indeksów
\(J\subseteq\{1,\ldots,m\}\)
wyznacza sympleks wtedy i tylko wtedy, gdy
\(\bigcap_{j\in J}U_j\neq\varnothing\).
Realizacja geometryczna nerwu
\(\lvert N(\mathcal U)\rvert\)
jest sumą odpowiednich ścian
standardowego sympleksu w \(\mathbb R^m\).
Jeśli rząd pokrycia jest co najwyżej \(n+1\),
realizacja nerwu ma wymiar symplicjalny
co najwyżej \(n\).
\end{definicja}

\begin{twierdzenie}[Odwzorowanie kanoniczne do nerwu]
\label{thm:topog23-canonical-nerve}
Niech \((X,d)\) będzie niepustą
zwartą przestrzenią metryczną,
a \(U_1,\ldots,U_m\) skończonym otwartym
pokryciem \(X\).
Istnieje ciągłe odwzorowanie
\(p\colon X\to\lvert N(\mathcal U)\rvert\)
o następującej własności:
dla każdego \(x\in X\) współrzędna
barycentryczna \(p_i(x)\) jest dodatnia
tylko wtedy, gdy \(x\in U_i\).
Jeśli rząd pokrycia jest co najwyżej \(n+1\),
obraz \(p(X)\) zawiera się w \(n\)-szkielecie nerwu.
\end{twierdzenie}
\begin{proof}
Dla każdego \(i\) określmy
ciągłą funkcję \(\psi_i\colon X\to[0,\infty)\)
w następujący sposób:
\[
 \psi_i(x)=
 \begin{cases}
   1,&U_i=X,\\
   d(x,X\setminus U_i),&U_i\neq X.
 \end{cases}
\]
Jeśli \(U_i\neq X\), dopełnienie
\(X\setminus U_i\) jest niepuste i domknięte,
więc funkcja odległości jest ciągła.
W obu przypadkach
\(\psi_i(x)>0\) dokładnie wtedy, gdy \(x\in U_i\).
Ponieważ \((U_i)\) pokrywa \(X\),
suma \(S(x)=\sum_i\psi_i(x)\) jest dodatnia
dla każdego \(x\in X\).
Przyjmijmy
\[
 p_i(x)=\frac{\psi_i(x)}{S(x)},
 \qquad
 p(x)=\sum_{i=1}^m p_i(x)e_i\in\mathbb R^m.
\]
Wszystkie współrzędne są ciągłe,
nieujemne i sumują się do \(1\).
Dodatnie współrzędne występują dokładnie
dla tych \(i\), dla których \(x\in U_i\).
Ich indeksy tworzą sympleks nerwu,
gdyż ich odpowiednie zbiory
mają wspólny punkt \(x\).
Stąd \(p(X)\subseteq\lvert N(\mathcal U)\rvert\).

Jeżeli rząd pokrycia nie przekracza \(n+1\),
dla każdego \(x\) co najwyżej \(n+1\)
współrzędnych \(p_i(x)\) jest dodatnich.
Tym samym \(p(x)\) leży na ścianie
o wymiarze nie większym niż \(n\).
\end{proof}

\begin{uwaga}[Związek z topologią algebraiczną]
\label{rem:topog23-nerve-overlap}
Nerw pokrycia występuje również
w topologii algebraicznej, gdzie
jego realizację bada się za pomocą
homologii, homotopii i twierdzenia o nerwie.
Tutaj potrzebujemy wyłącznie
konstrukcji odwzorowania barycentrycznego
i oszacowania liczby niezerowych współrzędnych.
Nie powtarzamy teorii homologii nerwu.
\end{uwaga}

\section{Górne oszacowanie wymiaru w przestrzeni euklidesowej}
\label{sec:topog23-euclidean}

\begin{lemat}[Drobna triangulacja przestrzeni euklidesowej]
\label{lem:topog23-grid-triangulation}
Niech \(n\ge1\) i \(h>0\).
Istnieje lokalnie skończona triangulacja
\(\mathbb R^n\), której wierzchołkami
są punkty siatki \(h\mathbb Z^n\),
a wszystkie sympleksy leżą w sześcianach
\[
 \prod_{j=1}^{n}[k_jh,(k_j+1)h],
 \qquad (k_1,\ldots,k_n)\in\mathbb Z^n.
\]
W tej triangulacji otwarte gwiazdy wierzchołków
tworzą pokrycie \(\mathbb R^n\)
rzędu co najwyżej \(n+1\).
Gwiazda wierzchołka \(v\)
zawiera się w \(v+(-h,h)^n\),
zatem ma średnicę nie większą niż \(2h\sqrt n\).
\end{lemat}
\begin{proof}
Rozważmy najpierw jednostkowy sześcian
\([0,1]^n\).
Dla każdej permutacji \(\sigma\)
zbioru \(\{1,\ldots,n\}\)
weźmy sympleks o wierzchołkach
\[
 0,\quad e_{\sigma(1)},\quad
 e_{\sigma(1)}+e_{\sigma(2)},\quad\ldots,\quad
 e_{\sigma(1)}+\cdots+e_{\sigma(n)}.
\]
Sympleks ten składa się z punktów,
których współrzędne w kolejności
\(\sigma(1),\ldots,\sigma(n)\)
są nierosnące.
Każdy punkt sześcianu spełnia
takie nierówności dla pewnej permutacji,
więc sympleksy pokrywają sześcian.
Dwa z nich przecinają się w ścianie
wyznaczonej przez równości
odpowiednich sąsiednich współrzędnych.
Te same podziały są zgodne na wspólnych
ścianach sąsiednich sześcianów.

Przesuwając i skalując konstrukcję
na sześciany siatki \(h\mathbb Z^n\),
otrzymujemy żądaną triangulację.
Jest lokalnie skończona, ponieważ
każdy zwarty zbiór spotyka tylko
skończenie wiele sześcianów siatki.
Dla wierzchołka \(v\)
jego \emph{otwarta gwiazda}
to suma wnętrz względnych sympleksów
mających \(v\) jako wierzchołek.
Równoważnie, jest to zbiór punktów,
dla których globalna, kawałkami afiniczna
współrzędna barycentryczna przy \(v\)
jest dodatnia. Współrzędna ta jest ciągła
dzięki zgodności triangulacji na ścianach,
więc gwiazda jest otwarta w \(\mathbb R^n\).
Każdy punkt należy do pewnej gwiazdy,
bo suma współrzędnych barycentrycznych
w odpowiednim sympleksie wynosi \(1\).

Dla punktu należącego do względnego wnętrza
sympleksu \(\sigma\)
dodatnie współrzędne barycentryczne
dotyczą wyłącznie jego wierzchołków.
Każdy sympleks ma co najwyżej \(n+1\)
wierzchołków, a więc punkt należy
do co najwyżej \(n+1\) otwartych gwiazd.
Każdy sympleks zawierający wierzchołek \(v\)
leży w jednym z sześcianów siatki
mających \(v\) jako wierzchołek.
Stąd gwiazda \(v\) zawiera się w
\(v+(-h,h)^n\), co daje oszacowanie średnicy.
\end{proof}

\begin{twierdzenie}[Górne ograniczenie wymiaru euklidesowego]
\label{thm:topog23-euclidean-upper}
Niech \(n\ge1\) oraz
\(K\subseteq\mathbb R^n\)
będzie zwartym podzbiorem.
Wówczas \(\dim K\le n\).
\end{twierdzenie}
\begin{proof}
Dla \(K=\varnothing\) teza wynika
z konwencji \(\dim\varnothing=-1\).
Dla niepustego \(K\)
weźmy dowolne skończone otwarte pokrycie
\((U_i)_{i=1}^m\) przestrzeni \(K\).
Niech \(\delta>0\) będzie jego liczbą Lebesgue'a
z Lematu~\ref{lem:topog23-lebesgue-number}.
Wybierzmy \(h>0\) takie, że
\(2h\sqrt n<\delta\),
i rozważmy triangulację z
Lematu~\ref{lem:topog23-grid-triangulation}.
Dla wierzchołków \(v\in h\mathbb Z^n\)
niech \(S_v\) będzie otwartą gwiazdą \(v\).

Rodzina \((K\cap S_v)\)
pokrywa \(K\) zbiorami otwartymi w \(K\)
i ma rząd co najwyżej \(n+1\).
Niepusty zbiór \(K\cap S_v\)
ma średnicę co najwyżej
\(2h\sqrt n<\delta\),
więc zawiera się w pewnym \(U_i\).
Ponieważ \(K\) jest ograniczony,
tylko skończenie wiele gwiazd spotyka \(K\):
jeśli \(S_v\cap K\neq\varnothing\),
to \(v\) leży w odległości
co najwyżej \(h\sqrt n\)
od pewnego punktu \(K\),
a w ustalonym ograniczonym
podzbiorze \(\mathbb R^n\)
jest tylko skończenie wiele punktów
siatki \(h\mathbb Z^n\).
Po odrzuceniu pustych przecięć
otrzymujemy skończone otwarte rozdrobnienie
wyjściowego pokrycia rzędu co najwyżej \(n+1\).
Zatem \(\dim K\le n\).
\end{proof}

\begin{uwaga}[Granice udowodnionych wyników]
\label{rem:topog23-euclidean-lower}
Powyższe twierdzenie dowodzi nierówności
\(\dim K\le n\), a nie równości dla wszystkich
zwartych \(K\subseteq\mathbb R^n\):
na przykład zbiór Cantora ma wymiar zero,
mimo że jest podzbiorem prostej.
Dla sześcianu \([0,1]^n\)
klasyczna równość \(\dim[0,1]^n=n\)
wymaga również dolnego oszacowania.
Wymaga ono dodatkowego twierdzenia
o pokryciach sześcianu lub równoważnego
argumentu topologicznego,
którego nie utożsamiamy
z udowodnioną tutaj nierównością górną.
Taki dowód powinien zostać umieszczony
w osobnej kontynuacji teorii wymiaru,
po sprawdzeniu odpowiednich wyników
w skrypcie topologii algebraicznej.
\end{uwaga}

\section{Obrazy ciągłe nie zachowują wymiaru}
\label{sec:topog23-images}

\begin{twierdzenie}[Obraz zbioru Cantora może mieć większy wymiar]
\label{thm:topog23-cantor-surjection}
Istnieje ciągłe odwzorowanie
\(f\colon C\to[0,1]\) zbioru Cantora
na cały odcinek.
W szczególności
\[
 \dim C=0,\qquad \dim f(C)=1.
\]
Zatem wymiar pokryciowy nie musi
maleć przy dowolnym ciągłym
odwzorowaniu surjektywnym.
\end{twierdzenie}
\begin{proof}
Z Przykładu~\ref{ex:topog23-cantor}
mamy homeomorfizm
\[
 c\colon\{0,1\}^{\mathbb N}\to C,\qquad
 c((\varepsilon_j))=\sum_{j=1}^{\infty}
                       2\varepsilon_j3^{-j}.
\]
Określmy odwzorowanie
\[
 b\colon\{0,1\}^{\mathbb N}\to[0,1],
 \qquad b((\varepsilon_j))=
             \sum_{j=1}^{\infty}\varepsilon_j2^{-j}.
\]
Jest ono ciągłe: suma pierwszych \(m\)
wyrazów zależy od skończenie wielu
współrzędnych, a ogon szeregu
ma moduł co najwyżej \(2^{-m}\),
jednostajnie względem wszystkich ciągów.
Jest także surjektywne,
ponieważ każda liczba z \([0,1]\)
ma co najmniej jedno rozwinięcie dwójkowe
o cyfrach \(0,1\); w przypadku \(1\)
wybieramy ciąg samych jedynek.
Wystarczy zatem przyjąć
\(f=b\circ c^{-1}\).
Ciągłość i surjektywność wynikają
z własności \(b\) oraz \(c\).
Równości wymiarów wynikają
z Przykładu~\ref{ex:topog23-cantor}
i Twierdzenia~\ref{thm:topog23-interval}.
\end{proof}

\section{Podsumowanie}
\label{sec:topog23-summary}

Wymiar pokryciowy mierzy możliwość
rozdrobnienia pokrycia otwartego tak,
aby w jednym punkcie spotykała się
ograniczona liczba jego elementów.
Dla zwartych przestrzeni Hausdorffa
wymiar zero jest równoważny istnieniu
bazy zbiorów otwarto-domkniętych.
W tym szczególnym wymiarze
otrzymaliśmy również zgodność
z małym i dużym wymiarem indukcyjnym.

Dowiedliśmy, że zbiór Cantora ma wymiar zero,
każdy nietrywialny zwarty zbiór
spójny ma wymiar co najmniej jeden,
odcinek ma wymiar jeden,
a każdy zwarty podzbiór \(\mathbb R^n\)
ma wymiar co najwyżej \(n\).
Nerw pokrycia wiąże ograniczenie rzędu
pokrycia z wymiarem kompleksu symplicjalnego.
Wskazaliśmy zarazem granicę obecnego
opracowania: niższe oszacowania wymiaru
sześcianów wyższych wymiarów
i pełne porównanie trzech wymiarów
wymagają dodatkowych wyników;
nie deklarujemy ich dowodów
na podstawie samych definicji.